\section{Adjustment vs.\ propagation, and sensitivity analysis}
\label{sec:propagation}

\emph{Source: handout \S1.5--1.6. Status: skeleton from a first read.}

\subsection{Adjustment vs.\ propagation (\S1.5)}

\newtag~study. Four arms compared: parents only; the efficient covariate set
(\S\ref{sec:estimation}); an estimator exploiting the DAG's factorization;
and full \textbf{parametric propagation} (fit every ECM structural equation,
simulate --- what the \emph{original} CDP paper's method does).

Findings to reproduce/understand, not just quote:
\begin{itemize}
  \item A confounder living entirely \textbf{downstream} (e.g.\ between
    weight and blood pressure, independent of exercise) does \emph{not}
    invalidate the parent-adjustment formula for the \emph{total} curve ---
    the formula is right whether or not that confounder is drawn in the
    graph. It changes the data (and so the curve), but adjustment still
    recovers the right curve for whatever the data is. Propagation only
    inherits this if its fitted conditionals match the observational ones
    --- which holds \emph{exactly} for linear OLS (an algebraic identity:
    fitted path coefficients always compose to the reduced-form slope) but
    can fail for nonlinear $f$: a Gaussian propagation model can match the
    mean and variance exactly and \emph{still} get $E[f]$ wrong (handout
    cites a checked example: error $0.047$ with mean and variance exact).
    \textbf{Nonlinearity is where propagation's extra assumptions can
    bite.}
  \item When the ECM's parametric equations are correct, propagation can be
    \emph{more precise} and can extrapolate outside the region of data
    overlap --- but that should be declared as model-based extrapolation,
    not treated as free. Overlap is diagnosed via the \textbf{non-overlap
    ratio} \reusetag{Bao \& Schomaker 2026}.
  \item Whether exploiting the DAG's factorization can beat plain adjustment
    at all is answered by \reusetag{Rotnitzky \& Smucler 2020, Thm.\ 19}.
\end{itemize}

\subsection{Sensitivity analysis (\S1.6)}

Mostly \reusetag{Cinelli \& Hazlett 2020; nonparametric bounds: Chernozhukov,
Cinelli, Newey, Sharma \& Syrgkanis}, \newtag~framing. If an unmeasured $U$
affects $A$ \emph{and} a downstream feature that matters for $f$, it can
bias the adjusted curve (special cases: the bias cancels, or $U$ turns out
irrelevant to $f$). Uses \textbf{partial-$R^2$ sensitivity} to ask: how
strong would $U$ need to be to flip the \textbf{slope} (never the level ---
that depends on an arbitrary zero point), flatten the curve, or cross a
threshold?

Project-specific (\newtag) content: confounding parameterized \emph{per ECM
edge}, with an exact map to the standard Cinelli--Hazlett coordinates (the
``tipping strength'' is their robustness-value formula, just
reparameterized); bounds at each bin; benchmarking against observed
predictors; and cross-checking against the \S1.5 invariance result above.

\begin{openquestions}
  \item Handout marks ``propagation arm via DoWhy $+$ invariance
    demonstrations'' as intern-feasible / moderate (\feasmod) --- a
    plausible task once \S\ref{sec:identification}/\S\ref{sec:estimation}
    basics are solid.
  \item ``Robustness section: linear closed forms, planted confounders,
    tipping, benchmarking'' is intern-feasible / moderate (\feasmod), with
    binned Riesz bounds ``only if time allows'' --- i.e.\ a stretch goal,
    not baseline scope.
  \item Work the toy linear-OLS ``path coefficients compose to the
    reduced-form slope'' identity by hand for the running example graph.
\end{openquestions}
